\documentclass[11pt]{article}
% ======================================================================
%  weekpreamble.tex — shared preamble for the weekly course summaries (EN)
%  Concise, readable, intuition-first. Same box style as the main doc.
%  Compiles with pdflatex.
% ======================================================================
\usepackage[utf8]{inputenc}
\usepackage[T1]{fontenc}
\usepackage[english]{babel}
\usepackage[a4paper,margin=2.2cm]{geometry}
\usepackage{amsmath,amssymb,amsthm,mathtools}
\usepackage{bm}
\usepackage{enumitem}
\usepackage{xcolor}
\usepackage[most]{tcolorbox}
\usepackage{microtype}

\emergencystretch=3em
\allowdisplaybreaks[1]
\setlength{\parindent}{0pt}
\setlength{\parskip}{0.45em}

% ---------- math macros (same names as the main document) ----------
\newcommand{\E}{\mathbb{E}}
\DeclareMathOperator{\Var}{Var}
\DeclareMathOperator{\Cov}{Cov}
\DeclareMathOperator{\Corr}{Corr}
\DeclareMathOperator*{\argmin}{arg\,min}
\DeclareMathOperator{\MSE}{MSE}
\DeclareMathOperator{\trace}{tr}
\newcommand{\Reals}{\mathbb{R}}
\newcommand{\Ints}{\mathbb{Z}}
\newcommand{\Comp}{\mathbb{C}}
\newcommand{\Nat}{\mathbb{N}}
\newcommand{\Prob}{\mathbb{P}}
\newcommand{\Tr}{^{\mathsf{T}}}
\newcommand{\Herm}{^{\mathsf{H}}}
\newcommand{\conj}[1]{\overline{#1}}
\newcommand{\abs}[1]{\left\lvert #1 \right\rvert}
\newcommand{\norm}[1]{\left\lVert #1 \right\rVert}
\newcommand{\ind}{\mathbf{1}}
\newcommand{\eps}{\varepsilon}
\newcommand{\diff}{\nabla}
\newcommand{\acvs}{\gamma}
\newcommand{\acf}{\rho}
\newcommand{\pacf}{\alpha}
\newcommand{\sdf}{\mathcal{S}}
\newcommand{\filt}{\mathcal{L}}
\newcommand{\Bsh}{B}
\newcommand{\ms}{\;\overset{\mathrm{ms}}{=}\;}
\newcommand{\toprob}{\xrightarrow{P}}
\newcommand{\todist}{\xrightarrow{d}}
\newcommand{\iid}{\overset{\text{iid}}{\sim}}

\setlist[itemize]{leftmargin=1.5em,topsep=2pt,itemsep=2pt}
\setlist[enumerate]{leftmargin=1.7em,topsep=2pt,itemsep=2pt}

% ---------- compact, titled (unnumbered) boxes ----------
\tcbset{wkbase/.style={enhanced,breakable,fonttitle=\bfseries,coltitle=white,
  boxrule=0.6pt,arc=1.2mm,left=2.2mm,right=2.2mm,top=1.1mm,bottom=1.1mm,
  before skip=6pt,after skip=6pt,
  attach boxed title to top left={yshift=-3.2mm,xshift=2.4mm},
  boxed title style={boxrule=0pt,arc=0.5mm}}}

\newtcolorbox{defn}[1]{wkbase, colback=blue!4!white, colframe=blue!58!black,
  colbacktitle=blue!58!black, title={Definition --- #1}, top=3mm}
\newtcolorbox{result}[1]{wkbase, colback=red!4!white, colframe=red!60!black,
  colbacktitle=red!60!black, title={#1}, top=3mm}
\newtcolorbox{intu}[1][Intuition]{wkbase, colback=yellow!12!white,
  colframe=orange!72!black, colbacktitle=orange!72!black, coltitle=white,
  title={#1}, top=3mm}
\newtcolorbox{retenir}[1][Key takeaways --- the week's reflexes]{wkbase,
  colback=green!5!white, colframe=green!48!black, colbacktitle=green!48!black,
  title={#1}, top=3mm}
\newtcolorbox{exmpl}[1][Worked example]{wkbase, colback=black!3!white,
  colframe=black!55!white, colbacktitle=black!55!white, title={#1}, top=3mm}

% ---------- week header ----------
\newcommand{\weekheader}[4]{%
  {\small\sffamily\bfseries\color{black!60} TIME SERIES~~$\cdot$~~MATH-342, EPFL}\par
  \vspace{2pt}
  {\LARGE\bfseries Week #1 --- #3\par}
  \vspace{1pt}
  {\itshape\color{black!55} #2}\par
  \vspace{6pt}
  \begin{tcolorbox}[enhanced,colback=blue!3!white,colframe=blue!45!black,
    arc=1.4mm,boxrule=0.7pt,left=3mm,right=3mm,top=2mm,bottom=2mm,before skip=2pt,after skip=10pt]
    \textbf{The big picture.}~ #4
  \end{tcolorbox}
}

\begin{document}
\weekheader{01}{19 February 2026}{Stationarity, Autocovariance \& Ergodicity}{Time series are dependent data; the whole analysis rests on first/second moments that are invariant to time shifts (weak stationarity) and on the autocovariance function.}

A \textbf{time series} $\{X_t\}$ is both the data recorded in time order \emph{and} the stochastic process $\{X_t:t\in T\}$ on $(\Omega,\mathcal F,\Prob)$ of which the data is one realisation. Here $T=\Ints$ and the sampling step $\Delta=1$. Unlike classical statistics, the observations are \textbf{dependent}: ``there are many more ways to be dependent than to be independent.'' We usually see only \emph{one} trajectory, yet must still recover population information from it.

\begin{intu}[The central challenge of the week]
With a single sample of length $n$, how can we estimate a ``population'' mean or variance? Answer: if the structure does not change over time (\textbf{stationarity}), then a \emph{time average} can stand in for an average over many realisations.
\end{intu}

\section*{Key definitions}

\begin{defn}{Weak (second-order) stationarity}
$\{X_t\}$ is \textbf{weakly stationary} (also: second-order, covariance stationary) if all first- and second-order moments exist, are finite, and are invariant to time shifts. Concretely, for all $t,s,\tau$:
\[
\text{(i) } \E[X_t]=\mu;\qquad
\text{(ii) } \Var(X_t)=\sigma^2<\infty;\qquad
\text{(iii) } \E[X_tX_{t+\tau}]=\E[X_sX_{s+\tau}].
\]
Condition (iii) forces $\E[X_tX_{t+\tau}]$ to be a function of the lag $\tau$ \emph{only}.
\end{defn}

\begin{defn}{Strict (strong) stationarity}
$\{X_t\}$ is \textbf{strictly stationary} if, for every $n$ and all $t_1,\dots,t_n$, the joint law of $(X_{t_1},\dots,X_{t_n})$ equals that of $(X_{t_1+\tau},\dots,X_{t_n+\tau})$: the \emph{whole} distribution is shift-invariant, not just the first two moments.
\end{defn}

\begin{intu}[Weak vs strict: neither implies the other]
\textbf{Strict $\not\Rightarrow$ weak}: i.i.d.\ Student-$t$ on $1$ d.o.f.\ (Cauchy) is strictly stationary but has no finite moments, so (ii) fails. \textbf{Weak $\not\Rightarrow$ strict}: weak stationarity says nothing about moments of order $\ge 3$. The two \emph{coincide} for a \textbf{Gaussian} process (every joint law is multivariate normal, hence fully fixed by mean and covariance).
\end{intu}

\begin{defn}{Autocovariance (ACVS, $\acvs_\tau$) and autocorrelation (ACF, $\acf_\tau$)}
For weakly stationary $\{X_t\}$,
\[
\acvs_\tau=\Cov(X_t,X_{t+\tau}),\qquad \acf_\tau=\frac{\acvs_\tau}{\acvs_0}=\Corr(X_t,X_{t+\tau}).
\]
Properties: $\acvs_0=\Var(X_t)\ge 0$, symmetry $\acvs_{-\tau}=\acvs_\tau$, and $\acf_0=1$, $\abs{\acf_\tau}\le 1$. The ACVS is, for a series, the analogue of the covariance matrix of a vector.
\end{defn}

\begin{defn}{ACVS estimators}
Given $X_1,\dots,X_n$ with sample mean $\bar X$, for $\abs\tau\le n-1$ (both $0$ otherwise):
\[
\tilde\acvs_\tau=\frac{1}{n-\abs\tau}\sum_{t=1}^{n-\abs\tau}(X_t-\bar X)(X_{t+\abs\tau}-\bar X),
\qquad
\hat\acvs_\tau=\frac{1}{n}\sum_{t=1}^{n-\abs\tau}(X_t-\bar X)(X_{t+\abs\tau}-\bar X).
\]
$\tilde\acvs_\tau$ (divisor $n-\abs\tau$) is \emph{unbiased} (with known mean) but need not be positive semi-definite. $\hat\acvs_\tau$ (divisor $n$) is \emph{biased} ($\E[\hat\acvs_\tau]=\tfrac{n-\abs\tau}{n}\acvs_\tau$) but \textbf{is} positive semi-definite: this is the default. Plug-in ACF: $\hat\acf_\tau=\hat\acvs_\tau/\hat\acvs_0$.
\end{defn}

\begin{defn}{Mean ergodicity}
$\{X_t\}$ is \textbf{mean ergodic} if its first two moments are finite and $\bar X\toprob \E[X_t]$ as $n\to\infty$: ``temporal averages $=$ population averages.'' Ergodicity and stationarity are \emph{not} equivalent.
\end{defn}

\section*{Key results \& formulas}

\begin{result}{Theorem --- Existence (Kolmogorov)}
A family of finite-dimensional distributions is that of a stochastic process \emph{iff} it is \textbf{consistent} (marginalising out one coordinate returns the law of the rest):
\[
\lim_{x_i\to\infty}F_{\mathbf t}(\mathbf x)=F_{\mathbf t(i)}(\mathbf x(i)).
\]
\textit{Why:} it guarantees that an infinite-index process $\{X_t\}_{t\in\Ints}$ really exists given only its finite marginals --- which legitimises modelling the whole time axis even though we observe only $n$ points.
\end{result}

\begin{result}{Proposition --- The ACVS is positive semi-definite}
For any $n$, all $t_1,\dots,t_n\in\Ints$ and all $a_1,\dots,a_n\in\Reals$:
\[
\sum_{j=1}^n\sum_{k=1}^n a_ja_k\,\acvs_{j-k}\ge 0 .
\]
\textit{Idea:} set $W=\sum_j a_jX_j$; then $0\le\Var(W)=\sum_{j,k}a_ja_k\acvs_{j-k}$. \textit{Why:} this is the \emph{defining} algebraic property of a genuine autocovariance --- a sequence that is not p.s.d.\ cannot be an ACVS.
\end{result}

\begin{result}{Proposition --- Variance of $\bar X$ and consistency}
If $\sum_\tau\abs{\acvs_\tau}<\infty$, then $\bar X$ is unbiased and
\[
\begin{aligned}
\Var(\bar X)&=\frac1{n}\sum_{\tau=-(n-1)}^{n-1}\Big(1-\tfrac{\abs\tau}{n}\Big)\acvs_\tau,
\qquad
n\,\Var(\bar X)\xrightarrow[n\to\infty]{}\sum_{\tau=-\infty}^{\infty}\acvs_\tau .
\end{aligned}
\]
Hence $\Var(\bar X)\to 0$ and $\bar X\toprob\mu$. \textit{Idea:} collect the double sum by lag $\tau$, then apply the \textbf{Ces\`aro summability theorem}. \textit{Why:} this is the prototype of ``when can a sample average replace a population average?''.
\end{result}

\begin{result}{Technique --- ACVS of a finite white-noise filter}
For $X_t=\sum_k\theta_k\eps_{t-k}$ with $\eps$ white noise of variance $\sigma^2$ (so $\Cov(\eps_s,\eps_u)=\sigma^2\ind_{\{s=u\}}$):
\[
\acvs_\tau=\sigma^2\sum_{j}\theta_j\,\theta_{j+\abs\tau}.
\]
Only the terms whose noise indices coincide survive --- this single reflex powers \emph{all} MA covariance computations.
\end{result}

\begin{intu}[Why dependence changes everything]
$n\,\Var(\bar X)\to\sum_\tau\acvs_\tau$: precision is driven not by $n$ alone but by the \emph{sum of all correlations}. Strong positive correlation (e.g.\ AR(1) with $\phi\to1$) inflates $\Var(\bar X)$: many observations, but little ``independent information.''
\end{intu}

\begin{intu}[Stationary $\ne$ ergodic]
Stationarity says the \emph{structure} does not move over time; ergodicity says \emph{one} trajectory suffices to recover the population mean. Typical counterexample: $X_t=Y$ with $Y$ a fixed random variable --- perfectly stationary, yet $\bar X\to Y$ (not a constant), so it is \emph{not} ergodic.
\end{intu}

\begin{exmpl}[Two-tap filter (MA(1))]
$X_t=\theta_1\eps_t+\theta_2\eps_{t-1}$ gives $\acvs_0=(\theta_1^2+\theta_2^2)\sigma^2$, $\acvs_{\pm1}=\theta_1\theta_2\sigma^2$, $\acvs_\tau=0$ otherwise, hence
\[
\acf_1=\frac{\theta_1\theta_2}{\theta_1^2+\theta_2^2},\qquad \abs{\acf_1}\le\tfrac12
\]
(by AM--GM, equality iff $\abs{\theta_1}=\abs{\theta_2}$). \textbf{No MA(1) can have $\abs{\acf_1}>\tfrac12$}: a handy sanity check.
\end{exmpl}

\section*{Key takeaways}

\begin{retenir}
\begin{itemize}
  \item \textbf{Test weak stationarity} in order, stop at the first failure: (1) is $\E[X_t]$ constant? (2) is $\Var(X_t)$ finite and constant? (3) does $\Cov(X_t,X_{t+\tau})$ depend on $\tau$ only? Classic trap: a covariance that keeps a $t$ (e.g.\ $\cos(ct)\cos(c(t+\tau))$).
  \item \textbf{Strict $\Leftrightarrow$ weak} only for \emph{Gaussian} processes; in general neither implies the other.
  \item \textbf{ACVS}: remember $\acvs_0=\Var\ge0$, $\acvs_{-\tau}=\acvs_\tau$, positive semi-definite. A non-p.s.d.\ sequence is not an ACVS.
  \item \textbf{Estimators}: by default divide by $n$ ($\hat\acvs_\tau$, biased but p.s.d.). Unbiasedness is not the goal; low MSE is.
  \item \textbf{Filtered white noise}: $\acvs_\tau=\sigma^2\sum_j\theta_j\theta_{j+\abs\tau}$ --- keep only the coinciding noise indices.
  \item \textbf{Consistency of $\bar X$} holds under $\sum_\tau\abs{\acvs_\tau}<\infty$; but dependence, not just $n$, controls precision. Stationary $\ne$ ergodic.
\end{itemize}
\end{retenir}

\end{document}
